\chapter{贾谊}

\begin{authorintro}
贾谊（前200—前168），西汉政治家、文学家。洛阳（今属河南）人。“年十八，以能诵诗属书闻于郡中。吴廷尉为河南守，闻其秀才，召置门下，甚幸爱”（《史记》本传）。因吴廷尉之荐，文帝召贾谊为博士，是时贾谊仅二十馀岁。每诏令议下，诸老先生不能言，贾谊尽为之对，人人各如其意所欲出。文帝悦之，一年之中超迁至太中大夫，欲任为公卿，引起周勃、灌婴等人的不满，文帝遂疏远贾谊，乃以贾谊为长沙王太傅，故称贾长沙。后迁为梁怀王（文帝少子）太傅，居数年，怀王骑马时堕马而死，贾谊自感未尽到为傅的责任，哭泣岁馀，亦死，时年三十三岁。有明张溥辑《贾长沙集》存世。今有王洲明、徐超校注本《贾谊集校注》（人民文学出版社1996年版）。
\end{authorintro}

\section{过秦论}

秦孝公\endnote{此文选自人民文学出版社《贾谊集校注》，但个别文字取用有异。《过秦论》是批评秦王朝的过失的，是贾谊《新书》中的名篇。汉初，鉴于秦的速亡，不少有识之士注意到总结历史上的经验与教训。刘邦初即位后，自认为天下是“居马上而得之”，曾骂爱称引《诗》、《书》的陆贾。陆贾反驳刘邦说：“居马上得之，宁可以马上治之乎？”刘邦因此改变了态度，“乃谓陆生曰：‘试为我著秦所以失天下，吾所以得之者何，乃古成败之国。’”陆贾著《新语》十二篇，每奏一篇，高祖未尝不称善，左右皆呼万岁。（事见《史记·郦生陆贾列传》）贾谊《新书》的产生背景及写作动机，与《新语》相类似。《过秦论》旨在通过分析秦王朝的兴亡（重点是灭亡）原因，劝谏汉文帝以秦为鉴戒，以求得汉帝国的长治久安。文章疏直峻切，气势雄伟，对比鲜明，且感情充沛，历来为人们所传诵。}据崤函之固\endnote{秦孝公（前381—前338）：秦献公之子，名渠梁，二十一岁即位，即位三年，用商鞅变法修刑，内务耕稼，外劝战死之赏罚，数年之后，秦国开始富强。崤（xiáo淆）函：崤，即崤山，在今河南洛宁西北，函，即函谷关，在今河南灵宝南，是秦的东关门户，东自崤山，西至潼津，深险如函，故名函谷。}，拥雍州之地\endnote{拥：拥有。雍州：古九州之一。《尚书·禹贡》：“黑水西河惟雍州。”在今陕西、甘肃、青海大部地区，为秦的封地所在。}，君臣固守，以窥周室\endnote{窥：窥测。周室：指东周王朝。}，有席卷天下、包举宇内、囊括四海之意\endnote{包举：统括；全部占有。囊括：包罗。义近包举。}，并吞八荒之心\endnote{八荒：八方极远之地。}。当是时也，商君佐之\endnote{商君：即商鞅，姓公孙，名鞅，战国时卫人，故又称卫鞅。由卫人秦后，为左庶长，劝秦孝公实行变法，秦国因之富强。后秦孝公封之于商，故号商君。}，内立法度，务耕织，修守战之具\endnote{修：治。守战之具：防守与出战的各种准备。具，一作“备”。}；外连衡而斗诸侯\endnote{连衡：即连横。东西为横，处于西方的秦与东方的齐、楚联合以打击其他国家叫连横。斗诸侯：使诸侯之间互相争斗。诸侯，指战国时的各国。}。于是秦人拱手而取西河之外\endnote{拱手：两手合在一起，表示毫不费力。西河：本战国魏地，今陕西东部黄河西岸地区。据《史记·商君列传》说，秦孝公使商鞅伐魏，魏使公子卬迎战，商鞅以与魏结盟为口实，诈骗公子卬，在会盟饮酒时，袭掠公子卬，又攻其军，魏大败，魏惠王乃派使者割河西之地献于秦，以求和。}。

孝公既没，惠文、武、昭襄王蒙故业\endnote{惠文：即秦惠文王，孝公之子，名驷，在位二十八年（前337—前310）而卒。武：即秦武王，惠文王之子，名荡，在位五年（前310—前306）而卒。昭襄：即秦昭襄王，乃秦武王的异母弟，名则，一名稷（据《史记·秦本纪》之《索隐》），在位五十六年（前306—前250）而卒。蒙故业：蒙受、继承秦孝公的旧有事业。}，因遗策\endnote{因遗策：遵循孝公留下的政策方略。}，南取汉中\endnote{南取汉中：据《史记·秦本纪》载：秦惠文王后元十三年（前311），“又攻楚汉中，取地六百里，置汉中郡”。}，西举巴、蜀\endnote{举：攻克；拔取。巴、蜀：古代国名，在今四川境内。《史记·秦本纪》载：秦惠文王后元九年，“司马错伐蜀，灭之”。}，东割膏腴之地，北收要害之郡\endnote{“东割”二句：《文选》李善注引李斯上书曰：“惠王用张仪之计，西并巴、蜀，南取汉中，东据成皋之险，割膏腴之壤。”膏腴之地，即肥沃的土地，当指韩国的宜阳（今属河南），魏的河东之地。要害之郡，当指成皋（今属河南）。但此地在秦之东，不在北，《史记》引此二句无“北”字，似较可靠。}。诸侯恐惧，同盟而谋弱秦\endnote{同盟为谋弱秦：指六国联合而谋划削弱秦国。同盟，《史记》所引作“会盟”。}，不爱珍器重宝、肥饶之地，以致天下之士\endnote{致：招纳。}，合从缔交，相与为一\endnote{“合从”二句：言六国采取合纵的策略，缔结同盟，互相援助，成为一体。南北曰纵，六国联合抗秦叫合纵。合从，同“合纵”。}。当此之时，齐有孟尝\endnote{孟尝：即孟尝君田文。《史记》有传。}，赵有平原\endnote{平原：即平原君赵胜。《史记》有传。}，楚有春申\endnote{春申：即春申君黄歇。《史记》有传。}，魏有信陵\endnote{信陵：即信陵君魏无忌。《史记》有传。此四人为战国时代的“四公子”，均蓄养大量的门客。}。此四君者，皆明智而忠信，宽厚而爱人，尊贤重士，约从离衡\endnote{约从离衡：缔约合纵，离散连横。}，兼韩、魏、燕、赵、宋、卫、中山之众\endnote{兼：一作“并”，又加；兼有。韩、魏、燕、赵、宋、卫、中山：皆战国时地处河南、河北等地的诸侯国，受秦威胁最大。}。于是六国之士\endnote{六国：指楚、齐、燕、韩、魏、赵。这是居于秦之东或东南的六大国，一时曾合纵而抗秦，但未见多少成效。}，有宁越、徐尚、苏秦、杜赫之属为之谋\endnote{宁越：《文选》李善注引《吕氏春秋》曰：“齐攻廪丘，赵使孔青将而救之，与齐人战，大败齐人，得尸三万，以为二京。宁越谓孔青曰：‘惜矣，不如归尸以内攻之，彼得尸而府库尽于葬，此之谓内攻之。’宁越，赵人也。”徐尚：未详。苏秦：东周洛阳人，《史记》有传。杜赫：《文选》李善注引《吕氏春秋》曰：“杜赫以安天下说周昭文君，昭文君谓杜赫曰：‘愿学所以安周。’高诱曰：‘杜赫，周人也。’”}，齐明、周最、陈轸、召滑、楼缓、翟景、苏厉、乐毅之属通其意\endnote{齐明：东周之臣。《文选》李善注引《战国策》：“东周齐明谓东周君曰：‘臣恐西周之与楚、韩宝（印信符玺），令之为己求地于东周也。’”周最：《文选》李善注引《战国策》曰：“齐令周最使郑，立韩扰而废公叔，周最患之。”高诱注：“周最，周君之子也，仕于齐，故齐使之也。”陈轸：《文选》李善注引《战国策》曰：“秦王谓陈轸曰：‘吾闻子欲去秦而之楚，信乎？’轸曰：‘然。’高诱注：‘陈轸，夏人，仕秦亦仕楚也。’”按：夏人，实为楚国人。夏，春秋时地名，楚庄王平陈国夏徵舒之乱，从陈每乡取一人聚居于此，称为夏州，见《左传·宣公十一年》。其地在今湖北汉阳北，战国时属楚国。召（shào绍）滑：一作“昭滑”，楚臣。《文选》李善注引《韩子》：“于象谓楚王曰：‘前时王使召滑之越，五年而能成之。’《史记》：‘范蠉对楚王曰：“王前尝用召滑而郡江东”。’召，音劭，滑音依字。”楼缓：《文选》李善注引《战国策》曰：“秦王伐楚，魏王不欲，楼缓谓魏王曰：‘不与秦攻楚，楚且与秦攻王，王不如令秦楚战，王交制之。’”高诱注：‘楼缓，魏相也。’”翟景：未详。苏厉：苏秦之弟，亦苏代之弟，先游说于燕，后奉燕质子入齐。齐王怨苏秦，欲囚苏厉。燕质子为谢，后苏代、苏厉皆归齐，而为齐臣。事见《史记·苏秦列传》。乐毅：本魏臣，使于燕，燕昭王以客礼待之，遂委质于燕，为燕之亚卿，懂兵法，为燕将，攻下齐七十二城。《史记》有传。通其意：沟通六国之士的意图。}，吴起、孙膑、带佗、倪良、王廖、田忌、廉颇、赵奢之朋制其兵\endnote{吴起：卫国人，闻魏文侯贤，事魏文侯以为将，是战国著名军事家之一，《史记》有传。孙膑：战国著名军事家，入齐后为齐威王的将军。事见《史记·孙子吴起列传》。带佗：楚将。倪良、王廖：皆为军事家。李善注引《吕氏春秋》：“王廖贵先（先法制人），兒良贵后（后法制人），此二人者，皆天下之豪士也。”兒，通“倪”。田忌：齐国大将。李善注引《战国策》：“韩、魏之君朝田侯，邹忌为齐相，田忌为将，使田忌伐魏，三战三胜。高诱注：‘田侯，（齐）宣王也。’”廉颇：赵国的名将。赵奢：赵国的大将，秦伐韩，赵惠文王曾令他将兵救韩，大破秦军。赐号马服君。之朋：之辈。《史记》作“之伦”，义同。}。尝以什倍之地，百万之众\endnote{百万之众：《史记》作“百万之师”。}，叩关而攻秦\endnote{叩关：叩击函谷关。叩关，一作“仰关”，此依《史记》。}，秦人开关延敌\endnote{延敌：将敌人引入。}，九国之师逡循而不敢进\endnote{九国之师：李善注：“九国，谓齐、楚、韩、魏、燕、赵、宋、卫、中山也。”逡循：同“逡巡”，犹豫不前。}。秦无亡矢遗镞之费\endnote{矢：指箭。镞：箭头。}，而天下诸侯已困矣\endnote{困：指处于困境。}。于是从散约解\endnote{从散约解：指六国“合纵”之约瓦解。}，争割地而赂秦。秦有馀力而制其弊\endnote{制其弊：利用六国的弱点而分别控制它们。}，追亡逐北\endnote{逐北：追击败北之敌。北，犹“败”。}，伏尸百万，流血漂橹\endnote{橹：大盾牌。一作“樐”。}，因利乘便\endnote{因利乘便：乘着有利与方便的时机。}，宰割天下，分裂山河。强国请伏\endnote{请伏：请求臣服。伏，通“服”。}，弱国入朝\endnote{入朝：指入秦朝拜秦王。}。

施及孝文王、庄襄王\endnote{施（yì义）：延续。孝文王：昭襄王之子，名柱，即位三天而卒。庄襄王：孝文王之子，名子楚，即位三年（前249—前247）而卒。}，享国日浅\endnote{浅：指时间短。}，国家无事。

及至始皇，奋六世之馀烈\endnote{奋：发扬。六世：指孝公、惠文王、武王、昭王、孝文王、庄襄王。馀烈：遗留的功业。}，振长策而御宇内\endnote{“振长策”句：以御马喻统治天下。振，挥动。长策，长鞭。御，驾驭，控制。宇内，天下。}，吞二周而亡诸侯\endnote{吞：吞并。二周：指西周、东周。西周建都洛阳西，东周建都巩（今河南巩义），分别在前256和前249年为秦所灭。亡诸侯：指灭六国等诸侯国。}，履至尊而制六合\endnote{履至尊：指登上帝位。履，足登其位。按：秦始皇于前246年继位，称秦王。二十六年后，始称始皇帝。六合：上下四方，实指普天之下。}，执搞朴以鞭笞天下\endnote{搞朴：同“敲朴”，棍棒。李善注引臣瓒曰：“短曰敲，长曰扑。”朴，亦作“扑”。鞭笞：鞭打，引申为统治。}，威振四海。南取百粤之地\endnote{百粤：一称“百越”，我国古代越族部落的总称，散居于今浙江、福建、广东、广西等地。越族支系繁多，除越国外，尚有瓯越、闽越、南越、骆越等，故称百越。此处之“百粤”，主要指广东、广西及越南北部地区。}，以为桂林、象郡\endnote{桂林：郡名，在今广西僮族自治区北部。象郡：郡名，在今广西南部，广东西南部，二郡为秦始皇时代新置。《史记·秦始皇本纪》载：“（始皇）三十三年，发诸尝逋亡人、赘婿、贾人略取陆梁（岭南）地，为桂林、象郡、南海，以适遣戍。”}；百粤之君，俯首系颈，委命下吏\endnote{委命下吏：委身听命于秦的下级官吏。}。乃使蒙恬北筑长城而守藩篱\endnote{蒙恬：秦始皇的将军。《史记·蒙恬列传》：“始皇二十六年，蒙恬因家世得为秦将，攻齐，大破之，拜为内史。秦已并天下，乃使蒙恬将三十万众北逐戎狄，收河南。筑长城，因地形，用制险塞，起临洮，至辽东，延袤万馀里。……是时蒙恬威振匈奴。”藩篱：本指篱笆，此代指边塞。}，却匈奴七百馀里\endnote{却匈奴：使匈奴退却。据《史记·匈奴列传》载，秦始皇使蒙恬将十万之众北击胡，悉收河南地。匈奴单于头曼，因不能胜秦而北徙。}，胡人不敢南下而牧马，士不敢弯弓而报怨\endnote{士：指匈奴的军士。弯弓：把弓张开，喻采取军事行动。报怨：报驱逐之仇。}。于是废先王之道\endnote{先王之道：指夏、商、周三代圣王的重礼义、仁政的治道。}，燔百家之言\endnote{燔（fán烦）：烧。百家之言：《诗》、《书》、诸子百家之言。焚书事见《史记·秦始皇本纪》。}，以愚黔首\endnote{以愚黔首：使百姓愚昧无知。秦代称百姓为黔首。}。堕名城\endnote{堕（huī灰）：毁坏。}，杀豪俊\endnote{豪俊：英雄豪杰。}，收天下之兵\endnote{兵：指兵器。}，聚之咸阳\endnote{咸阳：秦之都城，在今陕西咸阳东北。}，销锋镝\endnote{销锋镝（dì递）：熔化各种兵器。锋，指带刀刃的武器。镝，一作“镝”，箭头。}，铸以为金人十二\endnote{“铸以为”句：《史记·秦始皇本纪》：始皇二十六年，“收天下后，聚之咸阳，销以为钟鐻，金人十二，重各千石，置廷宫中”。}，以弱天下之民。然后践华为城\endnote{践华为城：据华山以为城郭。践，登，引申为据。}，因河为池\endnote{因河为池：依黄河为护城河。}，据亿丈之高\endnote{亿丈之高：形容华山为城之高。高，《史记》作“城”。}，临百尺之渊以为固\endnote{百尺之渊：指黄河为池之深。固：坚固可守。}。良将劲弩，守要害之处；信臣精卒\endnote{信臣：忠诚可靠之臣。}，陈利兵而谁何\endnote{谁何：呵问是谁，即盘问之义。}！天下已定，始皇之心，自以为关中之固\endnote{关中：秦以函谷关为门户，关中为秦的根据地，也是它的京畿地区。}，金城千里\endnote{金城：喻城郭的坚固如铁。}，子孙帝王万世之业也\endnote{帝王：作动词用，称帝称王。万世之业：《史记·秦始皇本纪》：“朕为始皇帝，后世以计数，二世三世至于万世，传之无穷。”}。

始皇既没，馀威振于殊俗\endnote{殊俗：指风俗不同于中原的边远少数民族地区。}。然而陈涉\endnote{陈涉：即陈胜，涉为其字。公元前209年，与吴广率戍卒九百人在大泽乡起义，震撼全国。}，瓮牖绳枢之子\endnote{瓮牖（yǒu有）：用破瓮作窗。绳枢：用绳子拴门轴。}，氓隶之人\endnote{氓隶：古代对农民的贱称。按：指“陈涉少时，尝与人傭耕”（《史记·陈涉世家》）。}，而迁徙之徒也\endnote{迁徙之徒：指被罚而服劳役的人。《陈涉世家》：“二世元年（前209年）七月，发闾左谪戍渔阳。”}。材能不及中人\endnote{中人：中等智力之人。}，非有仲尼、墨翟之贤\endnote{仲尼：孔子，字仲尼。墨翟：墨家创始人，著有《墨子》。}，陶朱、猗顿之富\endnote{陶朱：即陶朱公范蠡，他辅佐越王勾践灭吴后，以为勾践可以共患难，不可共富贵，辞官至陶（今山东定陶）经商，自号陶朱公。李善注引《史记》：“范蠡之陶，为朱公，以为陶天下之中，皆诸侯四通，货物所交易也，乃治产积十九年之间，三致千金。”猗（yǐ已）顿：鲁国的富翁。李善注引《孔丛子》：“猗顿，鲁之穷士也，耕则常饥，桑则常寒，闻朱公富，往而问术焉。公告之曰：‘子欲速富，当畜五牸（五种牲畜的雌种）。’乃适河东，大畜牛羊于猗氏之南，其滋息不可计。以兴富猗氏，故曰猗顿也。”}。蹑足行伍之间\endnote{蹑足：涉足；插脚。行伍：军队。此处指戍卒队伍。}，俛起阡陌之中\endnote{俛（fǔ府）起：俯仰。指劳作。阡陌：田间小路。此指农田。}，率疲弊之卒\endnote{疲弊：疲劳不堪。一作“疲散”。}，将数百之众，转而攻秦。斩木为兵\endnote{斩木：砍断木棍。兵：兵器，武器。}，揭竿为旗\endnote{揭竿：举竿。《文选》刘良注：“斩木为兵器而无锋刃，举竿为旗而无旌幡。”}，天下云合响应\endnote{云合：如云之集合。形容众多。响应：如响之应声，形容迅速。}，赢粮而景从\endnote{赢粮：担粮。李善注引《方言》：“赢，担也。”景从：如影之随形。景，同“影”。}，山东豪杰并起而亡秦族矣\endnote{山东豪杰：指项羽、刘邦等起义抗秦者。山东，指崤山之东。秦族：指嬴氏之族。}。

且夫天下非小弱也\endnote{天下非小弱：指秦之天下并未减少疆土、削弱力量。}，雍州之地、崤函之固自若也\endnote{自若：自如，和从前一样。}。陈涉之位，非尊于齐、楚、燕、赵、韩、魏、宋、卫、中山之君也；锄耰棘矜\endnote{锄耰（yōu幽）棘矜：代指陈涉起义时所用的武器。锄耰，即锄头。棘矜，带刺的木棍。}，不敌于钩戟长铩也\endnote{不敌：不能匹敌。一作“非铦”，铦（xiān先），锋利。钩戟：带钩的戟。长铩（shā杀）：大矛。钩戟长铩是当时作战的正式武器。}；谪戍之众\endnote{谪戍之众：指陈涉吴广所率领的九百戍卒。}，非抗于九国之师也\endnote{“非抗”句：不是九国军队的对手。}；深谋远虑，行军用兵之道，非及曩时之士也\endnote{曩时：过去，指战国时代。士：指战国时九国的军事家、政治家。}。然而成败异变\endnote{异变：异乎寻常的变化。}，功业相反也\endnote{功业相反：谓力量与功业成反比。}。试使山东之国与陈涉度长絜大\endnote{山东之国：指战国时代处于秦国东边的诸侯国。度长：比量长短。絜（xié协）大：比粗细。絜，用绳围量。}，比权量力，则不可同年而语矣。然秦以区区之地致万乘之势\endnote{区区：小小。致：达到，引申为获得。万乘之势：一作“万乘之权”，指天子的权势。周制：天子兵车万乘，诸侯兵车千乘。}，序八州而朝同列\endnote{序八州：安排天下州郡。意犹主宰天下。序，安排，序次。八州：古时天下分为九州，秦据雍州，秦以外还有八州。朝同列：使同列来朝拜。同列，指原先与秦地位平等的六国。}，百有馀年矣\endnote{百有馀年：秦统治天下不足二十年，说百有馀年不确，百，或为“十”之误。如从秦末上推百馀年，当在秦孝公、惠文王之世，但此时秦虽开始富强，并无“序八州而朝同列”之事。}。然后以六合为家\endnote{以六合为家：即以天下为家。喻秦搞中央集权，把天地四方看作嬴姓的天下。}，崤函为宫\endnote{崤函为宫：把崤山和函谷关以西的广大地区视作宫室。这是极力形容秦的私欲与占有欲。}。一夫作难而七庙堕\endnote{一夫作难：指陈涉一人首先发难反秦。七庙堕：七代祖先的宗庙被毁，实指秦王朝的灭亡。《礼记·王制》：“天子七庙。”}，身死人手\endnote{身死人手：指秦王子婴投降后为项羽所杀。}，为天下笑者，何也？仁义不施，而攻守之势异也\endnote{攻守之势异也：言以武力取天下和守成，形势是不同的。守成需施仁义，而秦用暴政，所以速亡。同时含告诫汉代皇帝，要广施仁义，天下能在马上得之，但不能马上治之。以上为上篇。}。

秦灭周祀\endnote{周祀：代指东周政权。按：古代新王朝建立，要封前二代的宗嗣为王，以不绝前代王朝的祭祀。秦始皇未遵古代礼制，不封周天子及六国诸侯的后代，灭周祀，不仅灭了东周政权，也断绝了周的祭祀。}，并海内，兼诸侯，南面称帝，以养四海\endnote{以养四海：为天下所供养。一本作“以四海养”。}。天下之士，斐然向风\endnote{斐然：犹“靡然”。顺风势而倒的样子。斐然向风，谓靡然景仰秦的风化。}。若是何也\endnote{若是何也：这是什么原因呢？}？曰：近古之无王者久矣\endnote{近古：指春秋战国。无王：周天子名存实亡，有王等于无王，故说无王。}。周室卑微\endnote{卑微：地位低下。指诸侯不把周天子看在眼里。}，五霸既灭\endnote{五霸：指春秋时先后称霸的齐桓公、晋文公、楚庄王、吴王阖闾和越王勾践。一说五霸指齐桓公、宋襄公、晋文公、秦穆公和楚庄王。五霸既灭，言到了战国时代。}，令不行于天下。是以诸侯力政\endnote{力政：以武力相征伐。政，通“征”。}，强凌弱，众暴寡\endnote{暴：施暴，欺侮。}，兵革不休，士民罢弊\endnote{罢弊：疲惫困乏。罢，通“疲”。}。今秦南面而王天下\endnote{王（wàng旺）天下：统治天下。王，作动词用。}，是上有天子也。即元元之民冀得安其性命\endnote{即：即便是。元元之民：老老实实的百姓。元元，忠实的样子。冀：希望。安其性命：保全他们的生命。}，莫不虚心而仰上\endnote{仰上：敬仰皇帝。}。当此之时，专威定功\endnote{专威：专有威权。指皇帝的至高无上的权威。定功：论定功勋以行封赏。实指下文的“裂地分民以封功臣之后”。}，安危之本，在于此矣。

秦王怀贪鄙之心\endnote{贪鄙：欲望大而见解偏狭。}，行自奋之智\endnote{自奋：自专骄矜。}，不信功臣，不亲士民，废王道而立私爱\endnote{王道：先王所行之正道。即以“仁义”治天下，与“霸道”相对。私爱：《史记》作“私权”，实指秦始皇实行的中央集权。}，焚文书而酷刑法\endnote{焚文书：焚烧书籍。酷刑法：指秦始皇所实行的严法峻刑。}，先诈力而后仁义\endnote{先诈力：重视诈术与武力。后仁义：轻视仁义道德。}，以暴虐为天下始\endnote{始：开端，创始。}。夫并兼者高诈力，安危者贵顺权\endnote{顺权：顺应形势。权，权宜。}，推此言之，取与守不同术也\endnote{术：方法。此句与上篇“攻守之势异也”意思相同。}。秦虽离战国而王天下\endnote{离战国：使战国时代的各国离散瓦解。离，一本作“并”，兼并。}，其道不易，其政不改，是其所以取之守之者异也\endnote{“是其所以”句：此句言秦取天下与守天下所采取的政策没什么不同。卢文弨本“也”前无“守之者异”四字，今据《史记》补。王念孙《读书杂志》指出：“异”上当有“无”字，否则文意不通，其说可从。}；孤独而有之\endnote{孤独而有之：独断专行地占有天下。孤独，指把所有的权力集中于皇帝一人。}，故其亡可立而待也。借使秦王论上世之事\endnote{借使：假使。论上世之事：指总结上古时代的治国之事。}，并殷周之迹\endnote{并殷周之迹：依傍殷周的足迹。并，通“傍”，依也。}，以制御其政\endnote{制御：控制。引申为治理。}，后虽有淫骄之主，犹未有倾危之患也。故三王之建天下\endnote{三王：夏、商、周三代开国之君，即夏禹王、商汤王、周文王。}，名号显美\endnote{显美：美名光显昭著。显，显扬。}，功业长久。

今秦二世立\endnote{秦二世：即秦始皇之子胡亥，秦始皇巡行，死于回来的路上，赵高、李斯互相勾结，矫诏立胡亥为帝，并迫使公子扶苏（始皇长子）自杀。}，天下莫不引领而观其政\endnote{引领：伸着脖子。表示急切的心情。观其政：观察二世的政治措施。}。夫寒者利裋褐而饥者甘糟糠\endnote{利裋（shù树）褐：以穿短粗布衣感到受益。裋，短衣。褐，粗布衣。甘糟糠：吃糟糠也感到甜美。}，天下嚣嚣\endnote{嚣嚣：喧哗、吵闹声。一本作“嗷嗷”。}，新主之资也\endnote{新主之资：这是新主胡亥给济百姓的机会。《庄子·大宗师》：“尧何以资汝。”注：“资者，给济之谓也。”}。此言劳民之易为政也\endnote{劳民：疲劳的百姓。易为政：容易治理。此句谓劳苦之民容易满足，故治理容易。}。向使二世有庸主之行而任忠贤\endnote{向使：那时假使。向，从前，表示过去的时间。庸主：平庸的君主。}，臣主一心而忧海内之患，缟素而正先帝之过\endnote{缟素：指丧服。代指丧期。正先帝之过：改正其父秦始皇的过错。}；裂地分民以封功臣之后\endnote{“裂地分民”句：分封土地与人民给功臣的后代。}，建国立君以礼天下\endnote{建国立君：封土建国以立六国之后代为君。以礼天下：以礼待天下之人。按：贾谊是主张分封（封建）制的。汉初，刘邦采纳郦食其的建议也曾想封六国的后代，但遭到张良的反对未能实行。贾谊此议并不高明。《东坡志林》卷五《秦废封建》条指出秦废封建，如“冬裘夏葛，时之所宜”，并明确指出：“封建者，争之端而乱之始也。”其说可参考。}；虚囹圄而免刑戮\endnote{虚囹圄（líng yǔ铃雨）：使监狱空虚。}，去收孥污秽之罪\endnote{收孥：收捕犯人的家属，即株连法。污秽之罪：当指黥、劓、宫刑等刑法。}，使各反其乡里；发仓廪\endnote{发：开。发仓廪，即开仓赈济。}，散财币，以振孤独穷困之士\endnote{振：通“赈”，救济。}；轻赋少事，以佐百姓之急\endnote{佐：帮助。}；约法省刑\endnote{约法：简化法令。省刑：减轻刑罚。}，以持其后\endnote{以持其后：以制约他们今后的行为。即以观后效之义。持，制约，挟制。《荀子·正名》：“以正道而辨奸，犹引绳以持曲直。”}，使天下之人皆得自新，更节循行\endnote{更节：改变操守。循行：遵循做人的正道。循行，《史记》作“修行”。}，各慎其身\endnote{各慎其身：各自珍重自身。即懂得小心谨慎的自我保护。}；塞万民之望\endnote{塞：堵塞。引申为杜绝。望：怨望；怨恨。}，而以盛德与天下\endnote{盛德：大德。与天下：使天下之人亲附。与，亲附；跟从。此为使动用法。}，天下息矣\endnote{息：安，太平。}。即四海之内，皆欢然各自安乐其处，惟恐有变，虽有狡害之民\endnote{狡害：《史记》作“狡猾”。}，无离上之心，则不轨之臣无以饰其智\endnote{饰其智：弄其巧，设其诈。《淮南子·本经》：“及伪之生也，饰智以惊愚，设诈以巧上。”}，而暴乱之奸弭矣\endnote{弭：止，消除。}。二世不行此术，而重以无道\endnote{重以无道：更加无道。}：坏宗庙与民\endnote{坏宗庙与民：毁灭宗庙与人民。按：《史记》之《集解》引徐广说：“一无此上五字”，认为此五字为衍文。}，更始作阿房之宫\endnote{更始：重新开始。防房之宫：即阿房宫，秦始皇三十五年始建，未建成而始皇死。二世元年四月，下令“复作阿房宫”（见《史记·秦始皇本纪》）。后被项羽烧毁。}；繁刑严诛，吏治刻深\endnote{刻深：犹“深刻”，指用法深而刻薄。据《秦始皇本纪》载：二世遵用赵高，申法令，恐大臣不服，诸公子与他争权，乃行诛大臣及诸公子，六公子戮死于杜，公子将闾无故被害，宗臣振恐，群臣谏者以为诽谤，大吏持禄取容，黔首振恐。此可作“繁刑严诛，吏治刻深”二句的注脚。}；赏罚不当，赋敛无度。天下多事，吏不能纪\endnote{纪：治理。}；百姓困穷，而主不收恤\endnote{收恤：收养抚恤。恤，同“恤”。}。然后奸伪并起\endnote{奸伪：奸诈和欺骗。}，而上下相遁\endnote{相遁：互相推脱以逃避责任。}；蒙罪者众，刑戮相望于道\endnote{刑戮相望于道：受刑被杀的人接连不断。}，而天下苦之。自群卿以下至于众庶\endnote{群卿：大臣。《史记》作“君卿”。众庶：众百姓。}，人怀自危之心，亲处穷苦之实，咸不安其位，故易动也\endnote{易动：容易发生变乱。}。是以陈涉不用汤武之贤，不借公侯之尊，奋臂于大泽\endnote{奋臂：振臂，挥动臂膀。大泽：即大泽乡（在今安徽宿州西南），为陈涉起义之处。}，而天下响应者，其民危也。

故先王者见终始之变\endnote{终始之变：事物发展变化的全过程。}，知存亡之由\endnote{由：原因。}。是以牧之以道\endnote{牧之以道：《史记》作“牧民之道”，即管理百姓的方法。牧，管理。道，道术，方法。}，务在安之而已矣\endnote{安之：使老百姓安定。《史记》此句无“矣”字。}。下虽有逆行之臣\endnote{下：《史记》作“天下”。}，必无响应之助。故曰：“安民可与为义，而危民易与为非\endnote{“安民”二句：言安定之民可以做正义之事，而身处危难之民则容易为非作歹。}”，此之谓也。贵为天子，富有四海\endnote{富有四海：拥有天下的财富。}，身在于戮者\endnote{身在于戮者：《史记》作“身不免于戮者”，于义较优。}，正之非也\endnote{正之非也：纠正败局的方法不对。《史记》作“正倾非也”。正倾，即纠正倾危之局势。}。是二世之过也\endnote{是二世之过也：这是二世胡亥的过错。按，以上为中篇。}。

秦兼诸侯山东三十馀郡\endnote{“秦兼”句：陶鸿庆《读诸子札记》说：此句疑有脱文，当云“秦兼山东诸侯，分天下为三十馀郡。”按《秦始皇本纪》云：“分天下以为三十六郡”，其中之上郡、陇西、巴郡、蜀郡等，非六国诸侯之地。陶说可从。}，循津关\endnote{循津关：依据渡口关塞。循，依，一本作“修”。}，据险塞，缮甲兵而守之\endnote{缮：修补，整治。甲兵：指武器装备。}。然陈涉率散乱之众数百，奋臂大呼，不用弓戟之兵，锄耰白梃\endnote{锄：同“锄”。耰：《史记》作“櫌”，锄柄。白梃：白色的木棒。}，望屋而食\endnote{望屋而食：行军不带粮食，走到哪里吃到哪里。}，横行天下。秦人阻险不守，关梁不闭，长戟不刺，强弩不射\endnote{“秦人”四句：言秦之险阻关塞与精良的武器，均未发挥丝毫的作用。按：此说系夸张，于史不合。二世元年七月陈涉起义后，九月项梁、项羽起义于会稽，随后刘邦亦起义，在楚师西进途中，秦军不断进行抵抗。章邯曾在定陶大败楚师，项梁战死。沛公、项羽去外黄攻陈留，陈留坚守不能下，足证秦兵曾顽抗。}。楚师深入\endnote{楚师：指陈涉的军队。涉起义后，称大楚、楚王，号张楚（张大楚国之义）。卢文弨本“楚师”作“楚沛”，误，据《史记》改。深入：指西进深入秦地。}，战于鸿门\endnote{鸿门：在今陕西临潼东。据《史记·陈涉世家》载，陈涉曾封周文为将军，命他西击秦。周文率车千乘，卒数十万，至戏，军焉。秦令大将章邯悉发兵以击楚大军，周文失败。章邯紧追周文，数月后周文自刭而死。戏，即戏亭，在鸿门附近，此言“战于鸿门”，似指这次战斗。}，曾无藩篱之难\endnote{曾无藩篱之难：言楚师深入时竟然没有感到突破秦的边防有何困难。藩篱，代指边防。}，于是山东诸侯并起，豪俊相立\endnote{豪俊相立：英雄豪杰相立为王。《史记》之《索隐》注豪俊为“武臣、田儋、魏豹之属。”又《秦始皇本纪》载：“（陈）胜自立为楚王，居陈，遣诸将徇地。山东郡县少年苦秦吏，皆杀其守尉令丞反，以应陈涉，相立为侯王，合从西乡，名为伐秦，不可胜数也。……武臣自立为赵王，魏咎为魏王，田儋为齐王。沛公起沛，项梁举兵会稽郡。”}。秦使章邯将而东征\endnote{章邯：原为少府，秦二世令其为大将，率骊山徒及家产奴隶和大军东征。先击破周文（周章）军，又出关杀陈胜于城父，破项梁于定陶，后投降了项羽。}，章邯因其三军之众要市于外\endnote{要市于外：章邯在钜鹿被项羽击败后，接受了陈馀的劝告，投降了项羽，与项羽“约共攻秦，分王其地”（《项羽本纪》）。要市，即约市，指互订契约来做交易。}，以谋其上\endnote{谋：图谋。上：指秦二世。}。群臣之不相信\endnote{不相信：《史记》作“不信”，指秦朝君臣之间互不信任。《秦始皇本纪》载：章邯数败之后，“二世使人让（责斥）邯，邯恐，使长史欣请事，赵高弗见，又弗信。欣恐，亡去，高使人捕追不及。欣见邯曰：‘赵高用事于中，将军有功亦诛，无功亦诛。’项羽急击秦军，虏王离，邯等遂以兵降诸侯。”}，可见于此矣。

子婴立\endnote{子婴：二世胡亥兄之子，二世被赵高杀死后，立子婴（又称公子婴）为秦王。为秦王四十六日于轵道投降刘邦。}，遂不悟\endnote{遂不悟：终竟不觉悟。遂，作副词用，终竟。}。借使子婴有庸主之材\endnote{借使：假使。庸主：平庸无能之主。}，而仅得中佐\endnote{中佐：中等的辅佐人才。}，山东虽乱，三秦之地可全而有\endnote{三秦：指秦国兼并六国之前的本土。项羽破秦入关后，三分秦关中之地，以封秦之降将，以章邯为雍王，领咸阳以西地；以司马欣为塞王，领咸阳以东至黄河地；以董翳为翟王，领上郡之地。合称三秦。}，宗庙之祀宜未绝也\endnote{宗庙之祀：宗庙的祭祀。周以前，新王朝建立，封前两代的后嗣为王，称作“二王”，以不绝前代帝王的宗庙之祀。自秦开始，废“二王”制度，前代王朝被推翻，宗庙之祀亦随之断绝。从此宗庙之祀能否继续，代表王朝的存亡。}。秦城被山带河以为固，四塞之国也\endnote{四塞之国：四面均有关塞的国家。}。自缪公以来\endnote{缪公：即秦缪公，名任好，公元前659年至前608年在位，他任用百里奚与蹇叔，晋惠公夷吾曾向他献河西之地，秦之国土遂东至黄河，他在发展秦国势力方面起过较大的作用。}，至于秦王，二十馀君，常为诸侯雄\endnote{诸侯雄：即称雄于诸侯。}，此岂世贤哉？其势居然也\endnote{“此岂世贤”二句：这难道是秦的世代之主都贤能吗？不，这是他们所处的地理形势优越才这样的。}。且天下尝同心并力攻秦矣\endnote{同心并力攻秦：指六国以合纵形式联合抗秦。}，然困于险阻而不能进者，岂勇力智慧不足哉？形不利，势不便也\endnote{“形不利”二句：言六国的形势不利。}。秦虽小邑，伐并大城，得阨塞而守之\endnote{“秦虽小邑”三句：言秦国虽然多为小城镇，却能征伐兼并大城市，这是因为秦国得以据险要之地且可以固守的原因。阨塞，险要关塞。}。诸侯起于匹夫\endnote{诸侯：指秦末起义的六国诸侯后代。因当时已无尺寸之权，又无一兵一卒，故言起于匹夫。}，以利会\endnote{以利会：因利益驱动而会合。}，非有素王之行也\endnote{素王：指有帝王之德而未有其位的人。}，其交未亲\endnote{其交未亲：指起义诸侯之间的关系并不密切。}，其下未附\endnote{其下未附：卢文弨本作“其名未附”，据《史记》改。附，归附，顺从。}，名曰亡秦，其实利之也\endnote{“名曰”二句：言起义诸侯名义上说是要推翻暴秦，其实是为各自的利益而战。}。彼见秦阻之难犯\endnote{彼：他们，指起义诸侯。秦阻：秦地险阻。}，必退师。安土息民\endnote{安土息民：安定本土使百姓得以休息。一作“安士息民”。}，以待其弊\endnote{以待其弊：以等待起义诸侯的疲困。弊，疲困。}，收弱扶罢\endnote{收弱扶罢（pí疲）：收容弱小的扶持疲困的。弱与罢，似指国内百姓。罢，通“疲”。}，以令大国之君\endnote{令：命令，使……服从。大国之君：指新建的东方诸侯国楚王、齐王、赵王等。}，不患不得意于海内。贵为天子，富有四海，而身为禽者\endnote{为禽：被擒。禽，同“擒”。此句指子婴在轵道投降刘邦。}，救败非也\endnote{救败非也：言子婴挽救败局的策略不对。}。

秦王足己而不问\endnote{秦王：指秦始皇。足己：自己认为满足，即骄傲自满之意。不问：不征求别人的意见。按：《史记·秦始皇本纪》说：“始皇为人天性刚戾自用，……意得欲从，以为自古莫及己。……乐以刑杀为威，天下畏罪持禄，莫敢尽忠。上不闻过而日骄，下慑伏谩欺以取容。”}，遂过而不变\endnote{遂过：遂顺错误，引申为坚持错误。不变：不改。}。二世受之\endnote{受之：指接受帝位。}，因而不改\endnote{因：指因循秦始皇的过错。}，暴虐以重祸\endnote{重祸：加重对人民及臣下的祸害。}。子婴孤立无亲，危弱无辅\endnote{危弱无辅：言子婴处于危急形势又无得力的辅佐大臣。}。三主之惑\endnote{三主：指秦始皇、二世胡亥与子婴。惑：迷惑。指迷惑于错误策略。}，终身不悟，亡不亦宜乎？当此时也，世非无深谋远虑知化之士也\endnote{知化之士：了解形势变化的人。}，然所以不敢尽忠拂过者\endnote{拂（bì毕）过：矫正错误。拂，通“弼”。}，秦俗多忌讳之禁也\endnote{忌讳之禁：主要指不得直言直谏。《史记·秦始皇本纪》载：秦始皇时，良士皆“不敢端言其过”，秦始皇“乃自除犯禁者四百六十馀人，皆阬之咸阳”。公子扶苏谏，不从，反而大怒，使扶苏到上郡去监蒙恬军。二世胡亥，指鹿为马，“群臣谏者，以为诽谤”。}。忠言未卒于口\endnote{未卒于口：话没说完。}，而身糜没矣\endnote{糜没：指被杀害。糜，通“靡”，烂。}。故使天下之士，倾耳而听，重足而立\endnote{重足而立：两只脚重叠起来，站立不前。形容小心谨慎。}，阖口而不言\endnote{阖口：合口，闭口。《史记》作“拑口”。}。是以三主失道，而忠臣不谏，智士不谋也\endnote{不谋：不出谋划策。}。天下已乱，奸不上闻\endnote{奸不上闻：坏事不让皇帝知道。奸，邪恶。}，岂不悲哉！

先王知壅蔽之伤国也\endnote{壅蔽：壅塞遮蔽。指下情不能上达，上情亦不能下达。}，故置公卿、大夫、士，以饰法设刑而天下治\endnote{饰法：整顿法度。设刑：建立刑法。}。其强也，禁暴诛乱而天下服，其弱也，五霸征而诸侯从\endnote{五霸：见本篇注〔113〕。此句言周室衰弱时，无力解决各国间的争端，由五霸来主持征伐事，代行天子之权，各诸侯国也服从。}，其削也，内守外附而社稷存\endnote{内守：在内部保住自己的国土。守，指守土。外附：对外依附各诸侯国的力量。社稷存：国家仍然存在，即保住了王室的名义。}。故秦之盛也，繁法严刑而天下震；及其衰也，百姓怨而海内叛矣。故周五序得其道\endnote{五序：卢文弨本作“王序”，据《史记》改。五序即周代的五等（公、侯、伯、子、男）序爵，代指周的分封制度。得其道：得其治道，得到治天下的法术。按：贾谊是主张封建制的，所以认为周代的五等爵制是得其道的。}，千馀载不绝\endnote{千馀载：周朝统治八百多年，言“千馀载”是夸张。}；秦本末并失\endnote{本末并失：先王以仁义为本，秦“仁义不施”，是失其本。乱已发生，“正倾”、“救败”措施不当，是失其末，故言本末并失。}，故不能长。由是观之，安危之统\endnote{安危之统：安定与危机的施政纲领。统，纲纪。}，相去远矣。

鄙谚曰\endnote{鄙谚：即俗谚。《史记》作“野谚”。}：“前事之不忘，后事之师也\endnote{“前事”二句：犹“前车之鉴”，义在总结历史的经验教训，以作为鉴戒与师法。}。”是以君子为国，观之上古，验之当世\endnote{验之当世：以上古治世之道实验于当代。}，参之人事\endnote{参之人事：参考于当世的人事变化。}，察盛衰之理，审权势之宜\endnote{审权势之宜：详察权势是否得当。}，去就有序\endnote{去就：取舍。有序：有一定的规则。有序则不乱。}，变化因时\endnote{因时：随时。}，故旷日长久而社稷安矣。

\printendnotes

\section{鵩鸟赋并序}

\begin{preface}
谊\endnote{此文选自《文选》卷十三，是贾谊谪居长沙时所作。《史记·屈原贾生列传》说：“贾生为长沙王太傅。三年，有鸮飞入贾生舍，止于坐隅。楚人命鸮曰鵩。贾生既已適（dí敌）居长沙，长沙卑湿。自以为寿不得长，伤悼之，乃为赋以自广。”这与《序》的内容大体相仿佛。此赋借同鵩鸟的对话，表达了作者的苦闷与悲观。其中有对人生的探索和哲理性的思考，这一切熔铸错综复杂的感情中，使作者陷入痛苦的深潭而不能自拔，最终不得不乞灵于《老》、《庄》之道，委身自然、乐天知命，以求得人生的解脱。刘勰《文心雕龙·诠赋》说：“贾生《鵩鸟》，致辩于情理。”此赋的艺术特点正是情理交融，开头以物兴感，又以感入理，以理见情，文笔劲健，格调亦复深沉。}为长沙王傅\endnote{长沙王：指刘邦所封的异姓王吴芮的玄孙吴差。见《汉书·吴芮传》。傅：即太傅，古代为三公之一。但诸侯国的太傅，地位低于皇帝的太傅，其地位次于诸侯国的国相。}，三年，有鵩鸟飞入谊舍\endnote{鵩（fú服）鸟：即猫头鹰。古人认为鵩为不祥之鸟。谚云：“不怕猫头鹰叫，就怕猫头鹰笑。”“猫头鹰上宅，不知死谁。”}，止于坐隅\endnote{坐隅：坐席的一角。}，鵩似鸮\endnote{鸮（xiāo消）：即猫头鹰。一说鸮与鵩是形貌相似的两种鸟。}，不祥鸟也。谊既以谪居长沙\endnote{谪（zhé哲）：贬官。}，长沙卑湿\endnote{卑湿：低洼潮湿。}，谊自伤悼\endnote{伤悼：忧伤。}，以为寿不得长，乃为赋以自广\endnote{自广：自我宽解。}。其辞曰：
\end{preface}

单阏之岁兮\endnote{单阏（chán è蝉饿）：卯年的别称。《尔雅·释天》：“太岁……在卯曰单阏。”《史记》之《集解》引徐广说：“文帝六年岁在丁卯。”按：北京大学中国文学史教研室选注《两汉文学史参考资料》云：“据清钱大昕《十驾斋养新录》及《廿二史考异》，则定丁卯为文帝七年，他说：‘徐氏不知古有超辰之法，故云六年也。’按，钱说是。因此可见贾谊初出为长沙王太傅时，当在文帝五年；及文帝七年，乃作《鵩鸟赋》，正合为太傅三年之数（参阅清朱珔《文选集释》）。”文帝七年，即前173年。}，四月孟夏\endnote{孟夏：夏季的第一个月。}。庚子日斜兮\endnote{庚子：查《中西回史日历》，是年四月庚子日为四月二十八日。日斜：太阳偏西之时。}，鵩集予舍。止于坐隅兮，貌甚闲暇\endnote{闲暇：不惊不慌、悠然自得。}。异物来萃兮\endnote{异物：指鵩鸟。萃：集，引申为止。一说萃，应作“椊”，止也（用王念孙《读书杂志》说）。}，私怪其故\endnote{私：暗自。故：缘故。}。发书占之兮\endnote{发：打开。书：指占卜一类的书籍。占：占卜，算卦。}，谶言其度\endnote{谶（chèn衬）：即谶语，预示吉凶的话。度：即数，定数。此句言书中的谶语已把吉凶的定数指示出来了。}。曰：野鸟入室兮，主人将去\endnote{去：离开此地。}。请问于鵩兮，予去何之？吉乎告我，凶言其灾\endnote{“吉乎”二句：言如果有好事，你就告诉我；如果有坏事，你也要把灾祸向我说明。}。淹速之度兮，语予其期\endnote{“淹速”二句：言自己的年寿是长是短，到底能活多久，请把期限告诉我。《文选》李善注：“淹，迟也；速，疾也。谓死生之迟疾也。”度，数。期，期限。}。鵩乃叹息，举首奋翼\endnote{奋翼：振动翅膀，将飞之貌。}，口不能言，请对以臆\endnote{请对以臆：《文选》李善注：“请以臆中之事以对也。”臆，胸臆，言鵩鸟心中明白，请示意作答。臆，《汉书》作“意”，即以贾谊之意会作为回答。}：

万物变化兮，固无休息。斡流而迁兮，或推而还\endnote{“斡（wò握）流”二句：言万物的运转推移，是循环反复的。斡流，运转。迁，移动。推，推移。还，回还往复。}。形气转续兮，变化而蟺\endnote{“形气”二句：言有形之物与无形之物互相转化连续不断，其变化就像蝉蜕皮一样。形，指有形体之物。气，指无形体之物。而，通“如”。蟺（chán蝉），通“蝉”。}。沕穆无穷兮\endnote{“沕（wù勿）穆”二句：言自然变化之理深微无穷，非言语所能说清。沕穆，深微，指上六句所言万物变化之理。}，胡可胜言。祸兮福所倚，福兮祸所伏\endnote{“祸兮”二句：语出老子《道德经》第五十八章。言祸福无常，往往福而祸生，而祸藏于福。倚，因，依。伏，藏，潜伏。}。忧喜聚门兮，吉凶同域\endnote{“忧喜”二句：言忧与喜、吉与凶常常是同聚一门，同在一处。域，处所。}。彼吴强大兮，夫差以败；越栖会稽兮，勾践霸世\endnote{“彼吴”四句：以春秋时代吴、越强弱胜败的变化，说明祸福吉凶的转换与无常。据《史记·越王句践世家》说：先是越王允常与吴王阖庐因互相征伐结下怨恨。允常死，其子句践继位，阖庐伐越，反被越王句践所败，阖庐被射伤，临死前，告其子夫差“必毋忘越”。夫差日夜练兵，准备报仇，句践闻之，先法治人，结果句践败于夫椒。句践派大夫种向吴王夫差请和。句践被赦反国后，栖息会稽，卧薪尝胆，发奋图强，终于灭了吴国，迫使夫差自杀。霸世，即为一代霸主。}。斯游遂成兮，卒被五刑\endnote{“斯游”二句：言李斯游说秦国得以成功，但最后还是受刑被害。李斯本楚国上蔡人，学成之后认为楚王不足事，六国皆弱，无可建功，遂西入秦，游说成功，位终宰相。二世二年七月，为赵高所谗，被腰斩咸阳市，夷三族。见《史记·李斯列传》。五刑，古代的五种刑罚。具体所指，其说不一。腰斩属五刑之一的大辟。}；傅说胥靡兮，乃相武丁\endnote{“傅说（yuè悦）”二句：言傅说虽身为刑徒，但终于做了殷高宗的宰相。傅说，本为奴隶，后为殷相。相传傅说曾执版筑于傅岩，武丁访得，与语大悦，举以为相，遂出现殷中兴的局面。《尚书·说命》、《吕氏春秋·求人》、《史记·殷本纪》均有记载。胥靡，是古代处分犯轻罪者的刑罚，用绳索将罪人系在一起，相随而行，以服劳役。胥即相义。靡，系也。武丁，殷高宗名武丁，他是殷代的中兴之主（见《史记·殷本纪》）。}。夫祸之与福兮，何异纠纆\endnote{“夫祸”二句：言祸与福相为表里，像两股绳与三股绳一样，常纠缠在一起。纠纆（mò墨），纠，为两股线撚成的绳。纆，为三股线撚成的绳子。}？命不可说兮，孰知其极\endnote{“命不”二句：言天命是不可用言语解说的，谁能知道它的终极呢？}？水激则旱兮，矢激则远\endnote{“水激”二句：言水受外力激发则流得快，箭受外力激发则射得远。激，指外力激发。旱，通“悍”，猛疾。}；万物回薄兮，振荡相转\endnote{“万物”二句：言万物互相逼迫激荡，引发种种变化。回，往反，来回。薄，逼迫。《易·说卦》：“山泽通气，雷风相薄。”回薄，往反逼迫，互相影响。振，同“震”。振荡，激荡。转，转化。《文选》李善注引《鹖冠子》曰：“水激则悍，矢激则远，精神迴薄，振荡相转。”为以上四句所本。}。云蒸雨降兮，纠错相纷\endnote{“云蒸”二句：以云雨等自然现象说明事物变化的错综复杂。蒸，气热而上升，蒸腾。纠错，纠缠错杂。纷，纷乱。}。大钧播物兮，坱圠无垠\endnote{“大钧”二句：言造化运转万物，使之运动发展，其范围是无边无际的。大钧，本为制陶者使用的转盘，此比作造化。播，转动。坱圠（yāng yà央亚），无边无际的样子。}。天不可预虑兮，道不可预谋\endnote{“天不可”二句：言天和道皆高深莫测，人的思虑与谋划对它是不能预先知道的。《文选》李善注引《鹖冠子》曰：“天不可预谋，道不可预虑。”}。迟速有命兮，焉识其时\endnote{“迟速”二句：言人的寿命长短自有天命决定，怎能知道具体时间呢？这两句是针对上文“淹速之度兮，语予其期”而发。}？

且夫天地为炉兮，造化为工；阴阳为炭兮，万物为铜\endnote{“且夫天地”四句：以冶炼金属为喻，说明万物的合散变化之道。炉，指冶炼金属的火炉。工，指冶炼金属的工匠。古代认为天地间的阴阳二气，是产生万物的母本，即阴阳所以成物，故把阴阳比作冶金的燃料“炭”。铜，为炼成之物，也由阴阳而成。}。合散消息兮，安有常则\endnote{“合散”二句：言万物的聚与散、生与灭，哪有一定的规律呢？合，聚合。散，分散。消，指灭。息，指生。}？千变万化兮，未始有极\endnote{未始：犹“未尝”。极：终极。}。忽然为人兮，何足控抟\endnote{“忽然”二句：言偶然生而为人，也不足珍贵，何必对它爱惜呢？忽然，犹偶然。控抟，持在手中抟弄。引申为爱惜珍重。}？化为异物兮\endnote{化为异物：指人死后变成另外的东西。}，又何足患？小智自私兮，贱彼贵我\endnote{“小智”二句：言眼光短浅的小智之人自私自利，以外物为贱，以己为贵。}。达人大观兮，物无不可\endnote{“达人”二句：谓通达之人见识远大，对万物一视同仁，故无所不宜。达人，指通达知命之人。大观，指所见远大。可，适宜。}。贪夫殉财兮，烈士殉名\endnote{“贪夫”二句：言贪财之人以身殉财，重义轻生之人，以身殉名。徇，同“殉”，以身从物曰殉。}。夸者死权兮，品庶每生\endnote{“夸者”二句：言好虚名、喜权势的人往往为追求权势而死，一般百姓则贪生而怕死。品庶，众庶，普通老百姓。每生，贪生。《文选》李善注引孟康曰：“每，贪也。”}。怵迫之徒兮，或趋东西\endnote{“怵迫”二句：言为利所诱、为贫穷所逼迫的人，总是东奔西跑地趋利避害。《文选》李善注引孟康曰：“怵，为利所诱怵也。迫，迫贫贱也。东西，趋利也。”}。大人不曲兮，意变齐同\endnote{“大人”二句：言大人不为外在的物欲所屈，故对千变万化，等量齐观。大人，指道德修养极高深的人。李善注：“大人者，与天地合其德。”曲，同“屈”，指为物欲所屈服。意，《史记》作“億”，億变，犹言千变万化。齐同，犹言等量齐观。实即庄子的“齐物论”，即把生与死、荣与辱看作一样。}。愚士系俗兮，窘若囚拘\endnote{“愚士”二句：言愚昧之人受世俗的羁绊，拘束得像被拘囚的犯人一样。系，牵，羁绊。窘，拘束。囚拘，犹拘囚。}。至人遗物兮，独与道俱\endnote{“至人”二句：言道德修养最高的人能遗世弃俗而不为外物所牵累，独能与大道同在。至人，道德修养达到最高境界的人，在大人之上。《庄子·逍遥游》：“至人无己，神人无功，圣人无名。”道，指老、庄所推崇的自然之道。}。众人惑惑兮，好恶积亿\endnote{“众人”二句：言一般人惑于世俗的利害，将好恶爱憎的感情积满胸中。惑惑，惑之甚也。亿，同“臆”，胸臆。}。真人恬漠兮，独与道息\endnote{“真人”二句：言得天地之道的人恬淡无为，所以独能与大道并存。真人，李善注引《文子》：“得天地之道，故谓之真人。”恬漠，淡泊无欲，虚静无为。息，生。引申为存在。}。释智遗形兮，超然自丧\endnote{“释智”二句：言弃绝智慧、忘掉自己的形体，超然于万物之外，嗒然自失。按：这是道家修养的境界，老子主张“绝圣弃智”，庄子的遗形、忘我，“心如死灰，形如槁木”与此二句思想一致。故李善注引《庄子》云：“仲尼问于颜回曰：‘何谓坐忘？’回曰‘堕支体，黜聪明，离形去智，同于大道，此谓坐忘。’”}。寥廓忽荒兮，与道翱翔\endnote{“寥廓”二句：言与混沌未分的元气同在，与大道合而为一。按：这是道家的修养境界。李善注：“寥廓忽荒，元气未分之貌。《广雅》曰：‘寥，深也；廓，空也。’《鹖冠子》曰：‘与道翱翔。’”}。乘流则逝兮，得坻则止\endnote{“乘流”二句：言人生如木浮水上，随流则飘浮而行，遇到阻碍就停止。按：这是道家的委命自然的思想。坻（chí池），水中小洲。}。纵躯委命兮，不私与己\endnote{“纵躯”二句：言放纵自己的躯体并把命运托付给自然，不把它看作是自己私有之物。纵，放纵。委，托付。李善注引《鹖冠子》曰：“纵躯委命，与时往来。”}。其生兮若浮，其死兮若休\endnote{“其生”二句：言活着好比把自己寄托在世上，死去就好比长久休息。浮，寄托。李善注引《庄子》曰：“其生如浮，其死若休。”}。澹乎若深泉之静，泛乎若不系之舟\endnote{“澹乎”二句：言人心的恬静应如无波的深渊那样宁静沉寂，浮游在世应像一只无绳拴住的舟船一样，任其自然地飘浮而无所沾滞。澹，恬静，安定。泛，漂浮。李善注引《庄子》：“老聃曰：‘其居也渊而静，其唯人心乎？’《鹖冠子》曰：‘泛泛乎若不系之舟。’”}。不以生故自宝兮，养空而浮\endnote{“不以”二句：言人不必因为活在世上的缘故就过于看重自己的生命，最好还是养其空虚之性，以浮游于人世。自宝，自珍，自贵。《汉书》注引服虔说：“道家养空虚若浮舟也。”}。德人无累，知命不忧\endnote{“德人”二句：言有修养的人是不会有所忧患的，因为他能知天命，所以不会有所忧愁。德人，有修养的人。累，忧患。李善注引淳芒曰：“德人者，居无思、行无虑也。”又曰“圣人循天之理，故无天灾，故无物累。”《周易·系辞上》曰：“乐天知命故不忧。”}。细故蒂芥，何足以疑\endnote{“细故”二句：言鵩鸟来舍，不过小事一桩，不必挂在心上。更不该有所疑虑。此句与开头几句呼应。细故，小事，琐碎的事故。蒂芥，即“芥蒂”。小的梗塞物，比喻心中有嫌隙或不适。}。

\printendnotes

\section{吊屈原文}

\begin{preface}
谊\endnote{贾谊的《吊屈原文》，一作《吊屈原赋》（见东汉王逸《楚辞注》）。本文选自《文选》卷六十。《文选》始改题为《吊屈原文》，并列为“吊文”类首篇。其实就它的文体形式看应视为“骚体赋”。《史记·屈原贾生列传》有“为赋以吊屈原”的话，并录其全文，为此赋的最早出处。此赋是贾谊被贬为长沙王太傅时渡湘水有感而作。长沙地方偏远，地势低洼潮湿，远谪长沙对贾谊来说是个很大的打击。他的遭遇与屈原的忠而见疑、遭谗放逐有相似之处。屈原的不幸遭遇，引起贾谊强烈的共鸣，吊屈原，亦自吊也。这是一首用骚体写成的抒情短赋，抒发的是遭贬后的愤懑之情，形式仿楚辞而趋于散文化，是骚赋向汉代大赋转型时代的代表作，在赋史上有承前启后的作用。}为长沙王太傅\endnote{长沙王：注见《鵩鸟赋》注〔2〕。}，既以谪去\endnote{谪（zhé哲）：被贬。去：指离开朝廷。}，意不自得\endnote{意不自得：即不得志。}，及渡湘水\endnote{湘水：即湘江。流经长沙。故汉时长沙称为临湘。}，为赋以吊屈原\endnote{为赋：即作赋。}。屈原，楚贤臣也，被谗放逐，作《离骚赋》，其终篇曰\endnote{终篇：篇末。}：“已矣哉！国无人兮，莫我知也\endnote{“已矣哉”三句：为《离骚》篇末之乱曰，原文为“已矣哉，国无人莫我知兮，又何怀乎故都”。已矣哉，犹言算了吧。莫我知，即莫知我。}。”遂自投汨罗而死\endnote{汨（mì密）罗：即汨罗江。汨水发源于江西修水，西南流入湖南，与发源于岳阳的罗水合流，故称汨罗江。}。谊追伤之，因自喻\endnote{自喻：比喻自己。意谓屈贾二人有相似的遭遇。“其辞曰”以上为后人所加之序文。}。其辞曰：
\end{preface}

恭承嘉惠兮\endnote{恭：敬。承：承蒙。嘉惠：好的恩惠。指皇恩。古代官吏被贬或流放，还要感谢皇帝的不杀之恩。}，俟罪长沙\endnote{俟罪：待罪。}。侧闻屈原兮\endnote{侧闻：从旁闻知。谦词。}，自沉汨罗。造托湘流兮\endnote{造托湘流：来到湘江边托流水来寄寓凭吊之情。造，至。托，寄托。}，敬吊先生。遭世罔极兮\endnote{遭世罔极：遭遇的时代无中正公平的准则。罔，无，不。极，中正的准则。}，乃殒厥身\endnote{殒（yǔn陨）厥身：丧失其身家性命。}。呜呼哀哉！逢时不祥\endnote{不祥：不好，不吉利。}。鸾凤伏窜兮\endnote{鸾凤：鸾鸟和凤凰，喻贤才。伏窜：逃避隐匿。}，鸱枭翱翔\endnote{鸱枭（chī xiāo吃消）：猫头鹰之类的猛禽。俗谓不祥之鸟，喻恶人。翱翔：飞舞。喻恶人横行。}。阘茸尊显兮\endnote{阘茸（tà róng沓荣）：庸碌低劣。尊显：尊贵显达。}，谗谀得志\endnote{谗谀：善进谗言、阿谀奉承的人。}。贤圣逆曳兮\endnote{逆曳：谓受迫而不能按正道行事。}，方正倒植\endnote{方正：品行端正之人。倒植：同“倒置”。}。世谓随夷为溷兮\endnote{随夷：指卞随与伯夷。皆为古代高尚的隐士。相传商汤将讨伐夏桀，曾和卞随商量，卞随拒不回答。商汤胜夏桀后，要让天下给卞随。卞随认为受到污辱，自投稠水（一说颍水）而死。事见《庄子·让王篇》、《吕氏春秋·离俗》等。伯夷，与其弟叔齐为殷末孤竹君之二子。周武王伐纣，二人扣马谏阻。武王灭商后，他们耻食周粟，采薇而食，后饿死首阳山。事见《吕氏春秋·诚廉》、《史记·伯夷叔齐列传》。溷：同“混”，污浊。指贪慕权势。}，谓跖蹻为廉\endnote{跖蹻（zhí jué侄决）：为古代著名的两个大盗。跖为鲁之大盗，名柳下跖。蹻为楚之大盗。廉：廉洁。}。莫邪为钝兮\endnote{莫邪：吴王阖庐使工匠铸造的宝剑。《荀子·性恶》：“阖闾之干将、莫邪、钜阙、辟闾，此皆古之良剑也。”钝：不锋利。}，铅刀为铦\endnote{铅刀：铅制的刀，铅质软，作刀不锐，为钝刀。铦（xiān先）：锋利。}。吁嗟默默\endnote{吁嗟：嗟叹。默默：不得意（用《汉书》颜师古注引应劭说）。}，生之无故兮\endnote{生：先生，指屈原。黄侃《文选平点》：“先生或省称先，或省称生。”无故：李善注引邓展曰：“言屈原无故遭此祸也。”一说无故即无罪。故，即“辜”的假借字。黄侃《文选平点》：“故即辜也。”两说皆可通。}！斡弃周鼎，宝康瓠兮\endnote{“斡（wò握）弃”二句：谓轻易抛弃了传国的宝鼎，却把破瓦器当作宝物。喻贤愚不辨，黑白混淆。斡，旋转。斡弃即抛弃。周鼎，周朝传国的宝鼎。宝，珍重，以……为宝。康瓠（hù户），空壶，破瓦壶。}。腾驾罢牛，骖蹇驴兮\endnote{“腾驾”二句：意谓想使车子飞奔却用疲惫的老牛驾辕，又用跛脚驴作骖马。喻用人不当。腾驾，使车子飞奔。罢（pí皮）牛，疲惫的老牛。骖（cān餐），在车辕两侧驾车的马叫骖。此作动词用。蹇（jiǎn简）驴，跛脚驴。}。骥垂两耳，服盐车兮\endnote{“骥垂”二句：意谓贤者不遇于时，大材小用。骥，千里马。垂两耳，无精打采的驯服之状。服，驾。盐车，载食盐的车。《战国策·楚策四》：“汗明曰：‘君亦闻骥乎？夫骥之齿至矣，服盐车而上太行。蹄申膝折，……白汗交流，中阪迁延，负辕不能上。……’”}。章甫荐履，渐不可久兮\endnote{“章甫”二句：言以帽子垫鞋，上下颠倒，如此倒行逆施，是难以持久的。章甫，殷代冠名。荐，垫。履，鞋子。渐，逐渐。}。嗟苦先生，独离此咎兮\endnote{“嗟苦”二句：对屈原的不幸遭遇表示同情。嗟苦，感叹劳苦。先生，指屈原。离，遭遇。咎，灾祸。}！

讯曰\endnote{讯曰：为结语，相当于《楚辞》中的“乱曰”。}：已矣！国其莫我知兮，独壹郁其谁语\endnote{壹郁：同“悒郁”、“抑郁”。忧郁、郁懑。谁语：犹“语谁”。向谁说，告诉谁。}？凤漂漂其高逝兮，固自引而远去\endnote{“凤漂漂”二句：言凤凰不愿与众鸟为伍而高飞远举，所以贤人亦应退隐而远去。凤，指凤凰。喻贤才。漂漂，同“飘飘”，高飞远举貌。高逝，远去。固，通“故”，所以。自引，主动退隐。}。袭九渊之神龙兮，沕深潜以自珍\endnote{“袭九渊”二句：言进入深渊的神龙，把自己深藏起来以保全自己不受伤害。袭，进入，入。九渊，深渊。沕（mì秘），潜藏。深潜，深藏。自珍，自我珍重，自我保全。}。偭猵獭以隐处兮，夫岂从虾与蛭螾\endnote{“偭猵獭（xiáo tǎ淆塔）”二句：言深藏的神龙背离水中的凶猛水族而隐居于深渊之中，又岂能追随小虾与蚯蚓呢？偭，违背。猵，鳄鱼类动物。獭，水獭。蛭（zhì致），水蛭，俗称蚂蟥。螾（yǐn引），蚯蚓。}？所贵圣人之神德兮，远浊世而自藏\endnote{“所贵”二句：言高洁之士，应当遗世而保持自己的品德。所贵，所珍贵的。神德，高洁的品德。远浊世，远离污浊的人世。}。使骐骥可得系而羁兮，岂云异夫犬羊\endnote{“使骐骥”二句：言假使千里马能被拴住和羁绊住，那与犬羊又有什么不同呢？骐骥，日行千里的骏马。系，拴住。羁，羁绊，束缚住。}。般纷纷其离此尤兮，亦夫子之故也\endnote{“般纷纷”二句：言在这乱纷纷的浊世中遭受祸患，也是由于夫子自身的原因。般，通“斑”，乱貌。纷纷，纷乱。离，同“罹”，遭受。尤，祸。夫子，指屈原。故，缘故。}。历九州而相其君兮，何必怀此都也\endnote{“历九州”二句：言屈原应当择主而仕，何必念念不忘楚国？历，经历，走遍。九州，代指全中国。相，观察。都，指楚都郢。}？凤凰翔于千仞兮，览德辉而下之\endnote{千仞：极言其高。八尺为一仞。德辉：道德的光辉。}。见细德之险征兮，遥曾击而去之\endnote{细德：指品德低微之人。险征：险恶的征兆。曾击：高飞。曾，高。击，展翅飞动。}。彼寻常之汙渎兮，岂能容夫吞舟之巨鱼\endnote{汙渎：亦作“污渎”。死水沟。汙，浊水不流。吞舟：能吞下舟船。极力形容鱼之巨大。}？横江湖之鳣鲸兮\endnote{横江湖：横断江湖之流水，形容形体之大。鳣（zhān沾）：一种大鱼，长可六七尺至二丈。}，固将制于蝼蚁\endnote{固：岂；难道。制：受牵制。蝼蚁：蝼蛄与蚂蚁。}！

\printendnotes